\section{Inverse RL: inferir la recompensa a partir del comportamiento}

\subsection{Idea básica}

Dadas las trayectorias de demostración del experto, se infiere qué \textbf{función de recompensa optimiza} el experto.

$$\max_\psi \; \mathbb{E}_{\tau \sim \pi_{\text{expert}}}\left[\sum_t r_\psi(s_t, a_t)\right] - \mathbb{E}_{\tau \sim \pi}\left[\sum_t r_\psi(s_t, a_t)\right]$$

Intuición: la recompensa aprendida debe hacer que el comportamiento del experto obtenga una recompensa más alta que el comportamiento de otras políticas.

\subsection{MaxEntIRL}

Al agregar la regularización de máxima entropía, IRL tiene una conexión profunda con las GAN: la función de recompensa es similar al discriminador y la política es similar al generador.

\subsection{Resumen de la sección}

Inverse RL infiere la función de recompensa a partir del comportamiento del experto, y es más flexible que el aprendizaje por imitación directa (puede manejar distintos embodiment), pero requiere optimizar de forma alternada la recompensa y la política.

